\documentclass[letterpaper,11pt]{article}
\usepackage[top=0.8in, bottom=0.8in, left=0.9in, right=0.9in]{geometry}
\usepackage[hidelinks]{hyperref}
\usepackage{lmodern}
\usepackage{parskip}
\usepackage{microtype}
\setlength{\parskip}{6pt}
\setlength{\parindent}{0pt}
\begin{document}
\begin{flushleft}
\textbf{\large Ajay Venkatesh} \\
Brooklyn, NY \\
+1 516 585 9013 \\
\href{mailto:av3855@nyu.edu}{av3855@nyu.edu}
\end{flushleft}
\vspace{10pt}
May 21, 2026
\vspace{6pt}

Hiring Manager \\
Stuller, Inc.

\vspace{10pt}

Dear Hiring Manager,

My name is Ajay Venkatesh. I'm finishing my M.S. in Computer Engineering at NYU, with production software experience at PwC in Bengaluru, a summer SDE internship at Amazon, and ongoing on-campus work at NYU's Novel AI Technologies. I'm writing about the Business Analyst role at Stuller.

The role's emphasis on writing and maintaining \textbf{SQL queries in Snowflake} to source, prepare, and validate data, then building \textbf{BI dashboards} (Strategy / MicroStrategy), maps directly onto the data-engineering work I've been doing. At PwC I spent the past few years on a distributed backend integrating with \textbf{Microsoft CRM} over \textbf{SQL Server}, writing complex queries, stored procedures, and parallelized batch ingestion that materially cut runtimes on operational feeds. I built \textbf{Power BI} dashboards over datasets joined across multiple relational databases for business stakeholders, exactly the source-prepare-validate-visualize loop the JD describes.

On ETL and large-scale data prep: my Big Data graduate project built a \textbf{PySpark / HDFS} pipeline over millions of flight records, with EDA, anomaly detection, feature engineering, and Random Forest / Gradient Boosted Tree models. At NYU's Novel AI Technologies I engineered \textbf{ETL pipelines} transforming \textbf{NoSQL} into \textbf{PostgreSQL} with concurrency control and locking mechanisms, then powered \textbf{AWS QuickSight} dashboards for real-time analytics. Identifying data issues, coordinating with technical partners, and validating outputs against source systems is the kind of careful work I've already been doing.

On stakeholder coordination and analytical project work: at PwC I owned features end-to-end alongside architects, product managers, and business stakeholders, gathering and documenting requirements and translating business questions into clear analytical work and reports. I'm comfortable with Jira-driven delivery, agile workflows, and the structured documentation that consulting-style analytics requires.

Stack-wise: \textbf{SQL}, \textbf{Python}, \textbf{Snowflake}, \textbf{PostgreSQL}, \textbf{MySQL}, \textbf{SQL Server}, plus \textbf{PySpark}, \textbf{Apache Spark}, \textbf{Pandas}, \textbf{NumPy}, and \textbf{Power BI}, \textbf{Tableau}, \textbf{QuickSight} for visualization. Comfortable across data warehousing, ETL, and dashboarding.

What pulls me toward Stuller is the operational depth of the work. Manufacturing and supply data for the jewelry industry runs on accuracy and clear reporting that real business users depend on every day, where careful, reproducible analytical work compounds across teams.

I'd welcome the chance to discuss further. Thanks for considering my application.

\vspace{12pt}
Sincerely, \\
Ajay Venkatesh

\end{document}